\section{Methods}
\label{sec:methodology}

\subsection{Speed Indicators}
\label{sec:meth_indicators}

For each trip $i$ with onboard speed vector $\mathbf{v}_i = (v_{i1}, \ldots, v_{iT_i})$ recorded at approximately 10-second intervals, we compute: mean speed $\bar{v}_i$, maximum speed $v_i^{\max}$, 85th percentile speed (P85), speed coefficient of variation ($\text{CV}_i = \sigma_i / \bar{v}_i$), and binary speeding indicator $\text{is\_speeding}_i = \mathbf{1}(v_i^{\max} > 25~\text{km/h})$ based on the South Korean regulatory limit. We also compute the speeding rate (fraction of speed readings exceeding 25~km/h), harsh acceleration count ($|a| > 0.5$~m/s$^2$, calibrated for the $\sim$10-second GPS interval), and cruise fraction (proportion of readings within $\pm$3~km/h of the median speed). User-level aggregates include speeding propensity (fraction of trips with any speeding), mode usage distribution, cumulative trip count, and experience measures.

\subsection{Cross-Sectional Model: Logistic GEE}
\label{sec:meth_gee}

We model binary speeding as a function of operating mode and controls via logistic regression with generalized estimating equations (GEE) to account for within-user correlation \citep{zeger1986longitudinal}:
\begin{equation}
    \text{logit}\big(P(\text{is\_speeding}_i = 1)\big) = \beta_0 + \beta_1 \text{TUB}_i + \beta_2 \text{ECO}_i + \boldsymbol{x}_i^{\top} \boldsymbol{\gamma}
    \label{eq:gee}
\end{equation}
where STD is the reference mode and $\boldsymbol{x}_i$ includes age group, time of day, day type (weekday/weekend), city fixed effects, fraction of trip on major roads, and log-transformed trip distance. We specify exchangeable working correlation and cluster by user ID. The key parameter of interest is $\exp(\beta_1)$, the odds ratio for speeding in TUB relative to STD mode. Robust (sandwich) standard errors are used throughout.

\subsection{Feature Importance: LightGBM with SHAP}
\label{sec:meth_shap}

To validate the parametric GEE findings and assess relative feature importance without distributional assumptions, we train a LightGBM gradient-boosted tree classifier on a stratified subsample of $\sim$2 million trips (80/20 train-test split, 5-fold cross-validation). The feature set includes all GEE covariates plus speed profile features (CV, cruise fraction, zero-speed fraction) and road class composition fractions (22 features total). Model performance is evaluated via AUC. SHAP (SHapley Additive exPlanations) values \citep{lundberg2017unified} are computed on the test set to quantify each feature's contribution to speeding predictions, providing a model-agnostic ranking of feature importance.

\subsection{Experience Analysis}
\label{sec:meth_experience}

To understand how speeding evolves with rider experience, we compute each user's cumulative trip rank across the February--November 2023 period and estimate a GEE logistic regression:
\begin{equation}
    \text{logit}\big(P(\text{is\_speeding}_i = 1)\big) = \beta_0 + \beta_1 \log(\text{trip\_rank}_i) + \beta_2 \text{TUB}_i + \boldsymbol{x}_i^{\top} \boldsymbol{\gamma}
    \label{eq:experience}
\end{equation}
clustered by user, with controls for mode, age, and temporal variables. The coefficient $\beta_1$ captures the experience--speeding gradient. To decompose this gradient, we separately analyze: (a) the relationship between experience and TUB mode adoption, (b) speeding trends within STD/ECO modes only, and (c) a mediation analysis quantifying the percentage of the experience effect attributable to mode adoption versus within-mode behavioral change.

\subsection{Difference-in-Differences: Causal Identification}
\label{sec:meth_did}

In December 2023, Swing banned TUB mode system-wide. Because cities varied in pre-ban TUB adoption, we exploit this variation in a two-way fixed effects (TWFE) DiD design. The panel comprises 53 cities $\times$ 11 months (February--December 2023). Our primary specification is:
\begin{equation}
    Y_{ct} = \alpha_c + \gamma_t + \beta \cdot (\text{Post}_t \times \text{TUB}^{\text{Nov}}_c) + \epsilon_{ct}
    \label{eq:twfe}
\end{equation}
where $Y_{ct}$ is the city-month speeding rate, $\alpha_c$ and $\gamma_t$ are city and month fixed effects, $\text{Post}_t = \mathbf{1}(t = \text{Dec 2023})$, and $\text{TUB}^{\text{Nov}}_c$ is the November 2023 TUB mode share in city $c$ (continuous treatment intensity). Observations are weighted by $\sqrt{N_{ct}}$ (square root of trip count) and standard errors are clustered by city \citep{arellano1987computing}. The parameter $\beta$ estimates the causal effect: a city with 50\% pre-ban TUB share is expected to experience a $0.5 \times \beta$ change in speeding rate.

We assess parallel trends via an event study specification:
\begin{equation}
    Y_{ct} = \alpha_c + \gamma_t + \sum_{k \neq \text{Nov}} \delta_k \cdot \mathbf{1}(t = k) \cdot \text{TUB}^{\text{Nov}}_c + \epsilon_{ct}
    \label{eq:event_study}
\end{equation}
with November as the reference period. Pre-period coefficients $\delta_k$ for $k < \text{Nov}$ are tested jointly ($F$-test) for parallel trends \citep{roth2022pretest}. We also estimate a dose-response relationship by regressing city-level speeding reduction (Nov $\rightarrow$ Dec) on TUB share reduction, with bootstrapped 95\% confidence intervals (10,000 replications).

\subsection{Robustness and Behavioral Compensation}
\label{sec:meth_robustness}

Three robustness checks address potential threats to causal validity. First, a \textbf{placebo test} applies the TWFE specification using September 2023 as a fictitious treatment date (restricting to February--October). A null result ($p > 0.05$) supports the identifying assumption. Second, \textbf{heterogeneity analysis} estimates the treatment effect separately for city-size quartiles. Third, a \textbf{de-trended event study} absorbs linear pre-trends to address the February--March ramp-up period.

To test for behavioral compensation, we estimate a \textbf{multi-outcome DiD} on five outcomes: (1) overall speeding rate, (2) trip distance, (3) STD/ECO-only speeding rate, (4) STD/ECO mean speed, and (5) STD/ECO maximum speed. If compensation occurs, outcomes 3--5 should show significant increases. We additionally conduct a \textbf{mode-switcher analysis}: among users who switched from TUB to STD/ECO after the ban (``switchers'') versus users who never used TUB (``never-TUB''), we compute within-user pre-post speed changes via a user-level DiD, reporting Cohen's $d$ effect sizes. A within-user placebo test using October 2023 as a fictitious ban confirms that any detected effects are treatment-driven.

\subsection{Survival Analysis: Time-to-First-Speeding}
\label{sec:meth_survival}

We model the time (in trips) to a user's first speeding event using survival analysis. The sample includes all users with $\geq$2 trips during February--November 2023. The event is the first trip where $v_i^{\max} > 25$~km/h; users who never speed are right-censored. Kaplan-Meier curves stratified by TUB usage (ever-TUB vs.\ never-TUB) are compared via the log-rank test. A Cox proportional hazards model estimates hazard ratios for TUB usage, age, and trip count. To address the endogeneity of TUB adoption timing, we additionally fit a Cox model with TUB adoption as a \textbf{time-varying covariate}: for each user-trip observation, TUB status indicates whether the user has adopted TUB by that trip. This captures the within-user effect of TUB adoption on speeding onset. Finally, we decompose the experience--speeding gradient by comparing hazard functions with and without TUB adoption to estimate the percentage of the experience effect mediated by mode choice.
